\documentclass[11pt,letterpaper,oneside]{report}
\usepackage[T1]{fontenc}
\usepackage[utf8]{inputenc}
\usepackage{lmodern}
\usepackage[margin=0.90in,headheight=15pt]{geometry}
\usepackage{microtype}
\usepackage[svgnames,table]{xcolor}
\usepackage{amsmath,amssymb,booktabs,longtable,tabularx,array}
\usepackage{enumitem}
\usepackage[most]{tcolorbox}
\usepackage{listings}
\usepackage{fancyhdr}
\usepackage{titlesec}
\usepackage{hyperref}
\usepackage{xurl}
\usepackage{lastpage}
\definecolor{ink}{HTML}{14273C}
\definecolor{primary}{HTML}{0B5379}
\definecolor{accent}{HTML}{B66A30}
\definecolor{pale}{HTML}{F1F5F8}
\definecolor{codebg}{HTML}{F5F7F8}
\hypersetup{colorlinks=true,linkcolor=primary,urlcolor=primary,pdftitle={Linux and UNIX System Programming Engineering Suite}}
\titleformat{\chapter}[display]{\normalfont\LARGE\bfseries\color{ink}}{\color{accent}\large\MakeUppercase{\chaptername}\ \thechapter}{10pt}{}[\vspace{0.4ex}\titlerule]
\titleformat{\section}{\Large\bfseries\color{primary}}{\thesection}{0.7em}{}
\titleformat{\subsection}{\large\bfseries\color{ink}}{\thesubsection}{0.7em}{}
\setlist[itemize]{leftmargin=*,topsep=3pt,itemsep=2pt}
\setlist[enumerate]{leftmargin=*,topsep=3pt,itemsep=3pt}
\setlength{\parskip}{5pt}
\setlength{\parindent}{0pt}
\lstdefinestyle{sysc}{language=C,basicstyle=\ttfamily\footnotesize,backgroundcolor=\color{codebg},keywordstyle=\bfseries\color{primary},commentstyle=\color{ForestGreen},stringstyle=\color{Brown},breaklines=true,showstringspaces=false,tabsize=4,frame=single,rulecolor=\color{LightSteelBlue},columns=fullflexible,keepspaces=true}
\lstset{style=sysc}
\newtcolorbox{engineeringnote}[1][]{colback=pale,colframe=primary,boxrule=0.55pt,arc=1.5mm,title=Engineering note,fonttitle=\bfseries,#1}
\newtcolorbox{modernnote}[1][]{colback=white,colframe=accent,boxrule=0.55pt,arc=1.5mm,title=Modern extension --- not in the 2010 textbook,fonttitle=\bfseries,#1}
\newcommand{\sys}[1]{\texttt{#1}}
\pagestyle{fancy}\fancyhf{}\fancyhead[L]{\small\color{ink} Linux / UNIX Engineering Suite}\fancyhead[R]{\small\color{primary}\nouppercase{\leftmark}}\fancyfoot[C]{\small\thepage\ / \pageref{LastPage}}
\renewcommand{\headrulewidth}{0.45pt}
\setcounter{tocdepth}{1}
\setlength{\emergencystretch}{2em}
\begin{document}
\begin{titlepage}\centering\vspace*{2cm}{\color{primary}\Large ENGINEERING REFERENCE SERIES\par}\vspace{1.25cm}{\color{ink}\Huge\bfseries Technical Reference\par}\vspace{0.6cm}{\Large Linux and UNIX system programming architecture\par}\vspace{2cm}\rule{0.75\textwidth}{1.2pt}\par\vspace{0.55cm}{\large An original companion to Michael Kerrisk\textquotesingle s 2010 text}\par\vspace{0.3cm}{\large Edition 1.0\quad |\quad October 2026}\par\vfill{\small Document 01\quad |\quad pdfLaTeX-compatible standalone source}\par\end{titlepage}
\tableofcontents\clearpage
\chapter{System programming mental model}
A user-space program interacts with the kernel through system calls, usually by C library wrappers. libc provides ABI mediation, cancellation behavior, feature-test-macro exposure, and additional functions implemented wholly in user space. Successful engineering begins by separating the API contract from kernel details, and persistent objects from process-local handles.

\section{Four separations to model}
\begin{itemize}
\item Process versus program: an exec operation replaces the program image without necessarily changing process identity.
\item File descriptor versus open file description versus inode: descriptor copies can share offsets; independent opens generally do not.
\item Virtual addresses versus physical memory: page tables, page faults, mappings and copy-on-write govern visibility.
\item Portability guarantee versus Linux-specific extension: the POSIX interface and Linux kernel facilities are not identical sets.
\end{itemize}
\begin{engineeringnote}
Document each API with caller-visible behavior, failure conditions, ownership transfer, synchronization constraints, resource bounds, privilege requirements, and portability. Do not infer behavior solely from a successful demonstration.
\end{engineeringnote}
\section{User-space / kernel boundary and ownership}
\begin{center}\fcolorbox{primary}{pale}{\parbox{.82\linewidth}{\centering
Application state, invariants and protocol framing\\[4pt]
\textbf{$\downarrow$ C/POSIX API and libc wrappers $\downarrow$}\\[4pt]
Kernel tasks, file tables, virtual memory, VFS, IPC and network stack\\[4pt]
\textbf{$\downarrow$ devices, peers, filesystems and scheduling $\downarrow$}\\[4pt]
Externally observable effects and recoverable failures}}
\end{center}
\section{A resource lifecycle contract}
\begin{enumerate}
\item Acquire the resource with explicit flags and a documented starting state.
\item Check the return contract; record only enough diagnostic context, then handle or propagate failure.
\item Establish a single responsible owner and define whether ownership is transferred or shared.
\item Perform partial and interrupted operations in a bounded loop; avoid consuming unvalidated data.
\item Apply deadlines, maximum sizes, and cancellation policies.
\item Close/unmap/destroy all resources; detect cleanup errors where they affect data integrity.
\item Assert externally observable postconditions and fault-injection behavior.
\end{enumerate}
\chapter{Portable interface design}
POSIX.1-2001 and POSIX.1-2008 frame the source textbook. POSIX.1-2024 is a later normative edition and must be verified independently. Linux-specific APIs should be hidden behind feature probes and explicit fallbacks where portability is a requirement.

\begin{longtable}{p{.22\linewidth}p{.31\linewidth}p{.35\linewidth}}\toprule
Boundary & Preferred baseline & Linux-only or implementation-sensitive alternative\\\midrule\endhead
File opening & \sys{openat}, \sys{O\_CLOEXEC} & \sys{openat2} with restricted resolution policies\\
Process launch & \sys{posix\_spawn}, \sys{fork}/\sys{execve} & \sys{clone3} and pidfds\\
Timing & \sys{clock\_gettime}, \sys{clock\_nanosleep} & \sys{timerfd} for descriptor-based loops\\
I/O readiness & \sys{poll} / \sys{select} & \sys{epoll}; \sys{io\_uring} for async completions\\
Filesystem watch & Periodic inspection with cost bounds & \sys{inotify}, fanotify (distinct capabilities)\\
Privileges & Effective credentials, permission model & Capabilities, seccomp, Landlock, namespaces\\
Resources & \sys{getrlimit}/\sys{setrlimit} & cgroup v2 hierarchical resource control\\
\bottomrule\end{longtable}
\section{Feature-detection protocol}
\begin{enumerate}
\item Build an API availability table for minimum kernel and libc, noting the syscall wrapper independently.
\item Probe runtime support without treating EPERM, ENOSYS and EINVAL as equivalent.
\item Keep an explicit contract for a supported fallback implementation.
\item Exercise unit and integration tests on the minimum supported systems; never gate solely on uname release strings.
\end{enumerate}
\chapter{I/O correctness and durability}
\section{Read/write semantics}
read may return fewer bytes than requested and zero indicates EOF for regular files; for sockets it can signal orderly stream shutdown. write may complete partially; robust loops track an offset and retry interrupted calls when appropriate.

\section{Atomicity and shared offsets}
Atomicity applies to specific interfaces and conditions, not arbitrary sequences. O\_APPEND controls append positioning; dup and fork can share an underlying open file description. pread/pwrite avoid the shared sequential offset for positioned operations.

\section{Durability and consistency}
stdio flushing, kernel page-cache writeback and power-loss persistence are distinct layers. For durable file replacement, a documented application may need temporary-file creation, data fsync, atomic rename and directory sync; filesystem and storage semantics still matter.

\section{Pathnames and identity}
Use dirfd-relative operations to bind traversal to known directories. Avoid check-then-open sequences on attacker-influenced paths. Retain file descriptors and verify the opened object with fstat when object properties matter.

\chapter{Process and thread lifecycle}
\section{Transition model}
\begin{itemize}
\item fork creates a child with copied logical memory state and inherited descriptor references; async-signal-safe restrictions matter between fork and exec in a multithreaded program.
\item execve replaces a process image. Deliberately specify argv/envp and close-on-exec policy; do not pass untrusted input through a shell.
\item waitpid and waitid reclaim child status; use a clearly owned reaping strategy to avoid zombies.
\item pthread\_create introduces concurrent execution over shared address space; pthread\_join coordinates completion and ownership.
\item Signals are event notifications with special safety rules, not a substitute for queue-like IPC or consistent shared-state synchronization.
\end{itemize}
\section{Thread/process invariant catalog}
\begin{itemize}
\item Every mutex protects a named invariant, and all accesses to the protected mutable state follow that convention.
\item A condition variable is paired with a predicate checked in a while loop.
\item All shared-memory protocols define initialization order, publication, atomicity and teardown.
\item Every spawned child has a reaping owner and an explicit termination path.
\item All waits have a stated policy for cancellation, signals, deadlines and shutdown.
\end{itemize}
\chapter{Filesystem and IPC decision architecture}
\begin{longtable}{p{.22\linewidth}p{.32\linewidth}p{.31\linewidth}}\toprule
Mechanism & Strong fit & Operational concern\\\midrule\endhead
Pipe / FIFO & Stream data between local processes & Backpressure, EOF and descriptor inheritance\\
UNIX stream socket & Bidirectional local protocol & Framing, authorization and peer identity\\
UNIX datagram socket & Bounded local messages & Delivery and size behavior must be defined\\
POSIX message queue & Prioritized messages & Kernel limits and persistent queue naming\\
Shared memory & High-throughput shared state & Synchronization and layout/ABI versioning\\
POSIX semaphore & Counting coordination & No automatic owner-based unlock\\
Record/file locks & Advisory coordinated file access & Lock namespace, ownership and crash cleanup\\
TCP socket & Reliable byte stream across hosts & No inherent message framing or TLS\\
UDP socket & Datagram exchange & Loss, ordering and size policies\\
\bottomrule\end{longtable}
\chapter{Modern Linux extensions and research map}
\section{openat2 and pathname resolution}
Use an open\_how structure and resolution controls for directory-scoped path resolution. The interface is Linux-only; do not mistake it for a general purpose security sandbox.

\begin{modernnote}
Read the latest platform documentation and test permissions, availability, and explicit fail-closed behavior before selecting this feature.
\end{modernnote}
\section{pidfd and clone3}
File-descriptor references to tasks support less racy lifecycle supervision. clone3 uses an extensible argument structure; version and kernel constraints require checking.

\begin{modernnote}
Read the latest platform documentation and test permissions, availability, and explicit fail-closed behavior before selecting this feature.
\end{modernnote}
\section{io\_uring}
A submission queue and a completion queue implement asynchronous operations. It is a completion model, not merely an interchangeable epoll optimization. Compare memory lifetimes, cancellation and supported opcodes.

\begin{modernnote}
Read the latest platform documentation and test permissions, availability, and explicit fail-closed behavior before selecting this feature.
\end{modernnote}
\section{seccomp, Landlock and no\_new\_privs}
seccomp filters restrict syscall entry points, not full object-level authorization. Landlock can restrict selected filesystem and other operations subject to ABI version and ruleset configuration.

\begin{modernnote}
Read the latest platform documentation and test permissions, availability, and explicit fail-closed behavior before selecting this feature.
\end{modernnote}
\section{Namespaces and cgroup v2}
Namespaces partition visibility and object identity; cgroups organize resource management and enforcement. Neither alone implies secure isolation.

\begin{modernnote}
Read the latest platform documentation and test permissions, availability, and explicit fail-closed behavior before selecting this feature.
\end{modernnote}
\section{eBPF and tracing}
eBPF enables controlled kernel instrumentation subject to verifier, privileges and deployment policy. Use tracing for diagnostics; do not claim that passive visibility prevents exploitation.

\begin{modernnote}
Read the latest platform documentation and test permissions, availability, and explicit fail-closed behavior before selecting this feature.
\end{modernnote}
\chapter{Engineering evidence and design review}
\begin{itemize}
\item Interface contract and version support matrix
\item Syscall trace capturing both successful and fault-induced paths
\item Thread ownership, locking order and race analysis
\item Input boundary and protocol framing tests
\item Linux privilege and sandbox configuration with least-permission justification
\item Latency distributions, resource-limit tests and operational cleanup observations
\end{itemize}

\chapter*{Source and currency notes}
\addcontentsline{toc}{chapter}{Source and currency notes}
\textbf{Foundational source:} Michael Kerrisk, \emph{The Linux Programming Interface: A Linux and UNIX System Programming Handbook}, No Starch Press, 2010, ISBN 978-1-59327-220-3. The textbook's chapter titles and start pages organize the study guide; explanations, tasks, and code in this suite are independently authored educational material, not reproductions of the original text or its example programs.

\textbf{Current reference baseline:} The Open Group / IEEE, \emph{POSIX.1-2024, Issue 8}, \url{https://pubs.opengroup.org/onlinepubs/9799919799/}; Linux man-pages project, \url{https://man7.org/linux/man-pages/}; Linux kernel documentation, \url{https://docs.kernel.org/}. For Linux-only extensions consult the relevant per-interface manual pages and kernel documentation for the target kernel, architecture and libc; availability can differ materially.

\textbf{Scope:} This is an engineering study and implementation companion. It does not claim to restate a normative standard or guarantee a specific Linux kernel's behavior. System calls, libc wrappers and higher-level library functions are explicitly distinguished where material.
\end{document}
